\def\writeANS{}
\begin{loigiaibth}{1}
  \begin {eqnarray*} \Delta _r H_{298}^\circ (1) &=& \Delta _f H_{298}^\circ (CaO) + \Delta _f H_{298}^\circ (CO_2) – \Delta _f H_{298}^\circ (CaCO_3)\\ &=& -635{,}1+(-393{,}5)-(-1206{,}9) = +178{,}3 \, kJ. \\ \Delta _r H_{298}^\circ (2) &=& \Delta _f H_{298}^\circ (Al_2O_3) + 2.\Delta _f H_{298}^\circ (Fe) – \Delta _f H_{298}^\circ (Fe_2O_3) – 2.\Delta _f H_{298}^\circ (Al)\\ &=& -1676 + 2.0 -(-825{,}5) - 2.0 = -850{,}5 \, kJ. \\ \Delta _r H_{298}^\circ (3) &=& \Delta _f H_{298}^\circ (CO_2) - \Delta _f H_{298}^\circ (CO) - \dfrac {1}{2} .\Delta _f H_{298}^\circ (O_2) \\ &=& -393{,}5 - (-110{,}5) - \dfrac {1}{2} .0 = -283 \, kJ. \\ \Delta _r H_{298}^\circ (4) &=& \Delta _f H_{298}^\circ (CO) + \Delta _f H_{298}^\circ (H_2) – \Delta _f H_{298}^\circ (C) – \Delta _f H_{298}^\circ (H_2O) \\ &=& -110{,}5 + 0 - 0 - (-285{,}84) = +175{,}34 \, kJ. \\ \Delta _r H_{298}^\circ (5) &=& 2.\Delta _f H_{298}^\circ (CO_2) + 3.\Delta _f H_{298}^\circ (H_2O) – \Delta _f H_{298}^\circ (C_2H_6) - \dfrac {7}{2} .\Delta _f H_{298}^\circ (O_2) \\ &=& 2(-393{,}5) + 3.(-285{,}84) - (-84) - \dfrac {7}{2} .0\\ &=& -1560{,}52 \, kJ. \end {eqnarray*}
\end{loigiaibth}
\def\writeANS{}
\begin{loigiaibth}{2}
 (a) \par \begin {tabular}{cccccccc} Phản ứng:&$C_6H_6(l)$ &+& $\dfrac {15}{2}O_2(g)$&$\longrightarrow $& $6CO_2(g)$& + &$3H_2O(g)$\\ $\Delta _fH_{298}^\circ $:&$+49{,}0$ && 0 &&$-393{,}50$ && $-241{,}82$\\ \end {tabular} \par $\Delta _rH_{298}^0 = 6 \times (-393{,}50) + 3 \times (-241{,}82) - 49 = -3135{,}46$ kJ \par \begin {tabular}{cccccccc} Phản ứng:& $C_3H_8 (g)$& +& $5O_2$ (g)& $\longrightarrow $& $3CO_2 (g)$ & + &$4H_2O(g)$\\ $\Delta _fH_{298}^\circ $:& $-105$ && $0$&& $-393{,}50 $&& $-241{,}82$\\ \end {tabular} \par $\Delta _rH_{298}^0 = 3 \times (-393{,}50) + 4 \times (-241{,}82) - 105 = -2042{,}78$ kJ \par (b) \par Số mol propane = $\dfrac {1}{44} \longrightarrow \Delta _rH_{298}^0 = \dfrac {-2042{,}78}{44} = -46{,}42$ kJ \par Số mol benzene = $\dfrac {1}{78} \longrightarrow \Delta _rH_{298}^0 = \dfrac {-3135{,}46}{78} = -40{,}2$ kJ \par Lượng nhiệt sinh ra khi đốt cháy $1{,}0$ g propane $C_3H_8(g)$ nhiều hơn so với lượng nhiệt sinh ra khi đốt cháy $1{,}0$ g benzene $C_6H_6(l)$.
\end{loigiaibth}
\def\writeANS{}
\begin{loigiaibth}{3}
  \begin {eqnarray*} \Delta _rH_{298}^\circ &=& 4\Delta _fH_{298}^\circ (CO_2) + 2\Delta _fH_{298}^\circ (H_2O) - 2\Delta _fH_{298}^\circ (C_2H_2) - 5\Delta _fH_{298}^\circ (O_2)\\ -2243.6&=& 4(-393.5) + 2(-285.8) - 2\Delta _fH_{298}^\circ (C_2H_2) - 5(0)\\ \Rightarrow \Delta _fH_{298}^\circ (C_2H_2) &=& +226.9 \text { kJ/mol} \end {eqnarray*}
\end{loigiaibth}
\def\writeANS{}
\begin{loigiaibth}{4}
  \begin {enumerate} \item Ba hydrocarbon X, Y, Z lần lượt là $\mathrm {HC \equiv CH}$ (ethyne hay acetylene); $H_2C=CH_2$ (ethene hay ethylene); $H_3C-CH_3$ (ethane). \item Phản ứng hoá học xảy ra: \\ $2C_2H_2(g) + 5O_2(g) \longrightarrow 4CO_2(g) + 2H_2O(g) \quad (3)$\\ $C_2H_4(g) + 3O_2(g) \longrightarrow 2CO_2(g) + 2H_2O(g) \quad (4)$\\ $2C_2H_6(g) + 7O_2(g) \longrightarrow 4CO_2(g) + 6H_2O(g) \quad (5)$ \item $\Delta _rH_{298}^0 (O_2) = 0$ \begin {eqnarray} \Delta _rH_{298}^0 (3) &=& 4 \times \Delta _fH_{298}^0 (CO_2) + 2 \times \Delta _fH_{298}^0 (H_2O) - 5 \times \Delta _fH_{298}^0 (O_2) - 2 \times \Delta _fH_{298}^0 (C_2H_2) \nonumber \\ &=& 4 \times (-393{,}5) + 2 \times (-241{,}82) - 2 \times (227{,}0) = -2\,511{,}64 \text { kJ} \nonumber \\ \Delta _rH_{298}^0 (4) &=& 2 \times \Delta _fH_{298}^0 (CO_2) + 2 \times \Delta _fH_{298}^0 (H_2O) - 3 \times \Delta _fH_{298}^0 (O_2) - \Delta _fH_{298}^0 (C_2H_4) \nonumber \\ &=& 2 \times (-393{,}5) + 2 \times (-241{,}82) - (52{,}47) = -1\,323{,}11 \text { kJ} \nonumber \\ \Delta _rH_{298}^0 (5) &=& 4 \times \Delta _fH_{298}^0 (CO_2) + 6 \times \Delta _fH_{298}^0 (H_2O) - 7 \times \Delta _fH_{298}^0 (O_2) - 2 \times \Delta _fH_{298}^0 (C_2H_6) \nonumber \\ &=& 4 \times (-393{,}5) + 6 \times (-241{,}82) - 2 \times (-84{,}67) = -2\,855{,}58 \text { kJ} \nonumber \end {eqnarray} \item Kết quả tính toán $\Delta _rH_{298}^0$ của phản ứng đốt cháy acetylene; ethylene; ethane giá trị lớn và $< 0$ (giải phóng năng lượng lớn) nên trong thực tiễn được sử dụng làm nhiên liệu. Riêng $C_2H_2$ trong thực tiễn làm đèn xi acetylene vì đèn xi acetylene có nhiệt độ cao nhất. \end {enumerate}
\end{loigiaibth}
\def\writeANS{}
\begin{loigiaibth}{5}
  Xét phản ứng: $2Al(s) + Fe_2O_3(s) \longrightarrow Al_2O_3(s) + 2Fe(s)$ \\ Biến thiên enthalpy của phản ứng: \begin {align*} \Delta _r H_{298}^\circ &= \Delta _f H_{298}^\circ (Al_2O_3) + 2.\Delta _f H_{298}^\circ (Fe) - 2.\Delta _f H_{298}^\circ (Al) - \Delta _f H_{298}^\circ (Fe_2O_3) \\ &= -1669{,}74 + 2.0 - 2.0 - (-822{,}4) = -847{,}34 \text { (kJ)} \end {align*} Nhiệt dung của sản phẩm: $C = 102{,}0{,}84 + 2.56.0{,}67 = 160{,}72 \text { (J.K}^{-1}\text {)}$. \\ Nhiệt độ tăng lên: $\Delta T = \dfrac {847{,}34.10^3.50 \%}{160{,}72} = 2636 \text { (K)}$ \\ Nhiệt độ đạt được: $(25 + 273) + 2636 = 2934 \text { (K)}$
\end{loigiaibth}
\def\writeANS{}
\begin{loigiaibth}{6}
  Phản ứng cháy: $C_2H_5OH(l) + 3O_2(g) \longrightarrow 2CO_2(g) + 3H_2O(l)$ \begin {align*} \Delta _r H_{298}^\circ = 2.(-393{,}5) + 3.(-285{,}8) - (-267) = -1377{,}4 \text { (kJ)} \end {align*} Nhiệt lượng cần cung cấp để đun $100$ gam nước từ $25^\circ C$ đến $100^\circ C$ là $Q = 100.4{,}2.(100 - 25) = 31500 \text { J} = 31{,}5 \text { kJ}$ \\ Gọi khối lượng cồn cần đốt là m (g) $\Rightarrow \dfrac {m}{46}.1377{,}4.60 \% = 31{,}5 \Rightarrow m = 1{,}75 g$.
\end{loigiaibth}
